\section{Socioeconomic impacts}

\section{Economic impacts}

The economic impacts of \acp{bss} are multi-layered, and while most \acp{bss} operate with public support or sponsorship and typically do not generate direct profits~\cite{Luo2019}, their economic relevance becomes more apparent once externalities such as travel time savings, health improvements, and local spending are considered.

At the user level, bike-sharing can reduce both monetary and time costs, especially for users shifting from car, ride-hailing, or \ac{pt}. However, because many trips replace walking or \ac{pt} rather than car use~\cite{Murphy2015, Fishman2013, Fishman2014, BachandMarleau2012}, direct per-trip financial savings are often limited, whereas time savings can still be meaningful given the economic value typically assigned to travel time savings~\cite{Qiu2018}. Aggregated across users, these effects can be substantial. In Beijing, bike-sharing as a last-mile commuting solution was estimated to save about 8 minutes per worker per day, yielding annual productivity gains of 0.3--1.2 billion yuan depending on assumptions~\cite{Qiu2018}. In Dublin, \textit{Dublinbikes} was associated with annual savings of 16.5 million minutes compared to prior trips (often walking or \ac{pt}), valued at about 6~million~\texteuro{} in direct time benefits, plus an additional 6.79~million~\texteuro{} per year in wider agglomeration benefits through increased accessibility and effective density~\cite{Bullock2017}. Evidence from London similarly indicates \ac{bss} trips are, on average, about 20\% shorter than the journeys they replace~\cite{Woodcock2014}, and U.S. evidence suggests bike-sharing can be as fast or faster than taxis for short \ac{od} trips (under 3~km) during rush-hour congestion~\cite{FaghihImani2017_Hail}. Overall, these findings suggest that bike-sharing can contribute to more time-efficient urban mobility.

Beyond travel-time effects, bike-sharing can stimulate local commerce by improving access and enabling additional trips. Survey evidence from \textit{Nice Ride Minnesota} suggests users made extra trips--primarily to retail and dining destinations--because of bike-sharing convenience, reporting nearly \$35,000 in expenditures in one week~\cite{Schoner2012}. In Washington, D.C., \textit{Capital Bikeshare} users reported making trips they would not otherwise have taken (16\%) and spending more at destinations reached by bike (23\%), while businesses near stations reported increased sales and positive neighborhood impacts~\cite{Buehler2014}. Although some effects may reflect redistribution rather than net new spending, the evidence supports bike-sharing as a contributor to local commercial activity and visibility. Moreover, bike-sharing has also been linked to property-market effects. In Montreal, proximity to stations was associated with higher multifamily housing sale prices. Each additional station within 800 meters corresponded to about a \$709 increase in unit price, implying roughly \$8,650 per unit (about 2.7\% of average prices) given typical station densities in the study area~\cite{ElGeneidy2016}.

Health benefits add a further economic dimension. Across 12 European cities, annual health benefits attributable to \acp{bss} use have been estimated at 18 million~\texteuro{}~\cite{Otero2018}, while in the U.S. annual health-related economic gains have been valued at around \$36 million~\cite{Clockston2021}. A cost--benefit analysis of \textit{Dublinbikes} found that benefits--notably health-related--outstripped costs, with benefit-cost ratios ranging from 1.21 to 1.68, ensuring a net societal gain for every euro invested~\cite{Murphy2015}.

Despite these diverse positive impacts, potential negative economic consequences should not be overlooked. The rapid expansion of dockless \acp{bss} in China around 2017 illustrates such risks: some major operators experienced mass layoffs and withdrew from multiple cities~\cite{Fortune2017, BusinessInsider2018}, and public agencies eventually bore costs related to abandoned bikes and urban clutter~\cite{Chen2019}. Thus, as cities leverage the economic upsides of bike-sharing, prudent planning and regulation remain important to mitigate wider costs.

\section{Social equality}

While \acp{bss} are frequently promoted as instruments for advancing social equity~\cite{Shaheen2010}, their ability to do so depends heavily on who can access and use them. Because proximity to stations strongly predicts usage~\cite{Kabra2018, BachandMarleau2012}, unequal spatial coverage can reinforce existing mobility inequalities, particularly among disadvantaged groups affected by transport poverty (e.g., low car ownership and limited \ac{pt} access)~\cite{Lucas2016}. Accordingly, research consistently identifies spatial distribution as a key driver of inequality in bike-sharing~\cite{Teixeira2021, Goodman2014, Ogilvie2012}. Stations are often concentrated in central, dense, and affluent neighborhoods, leaving peripheral or lower-income areas with weaker access, a pattern documented in cross-city spatial fairness analyses~\cite{DuranRodas2020}. Although these strategies maximize ridership by targeting major activity hubs, they may also contribute to persistent socio-demographic imbalances in the user base (see Section~\ref{sec:who_uses_bss})~\cite{Ogilvie2012, Woodcock2014, Ricci2015}.

These disparities are increasingly discussed within broader frameworks of transport equity and distributive justice~\cite{Cook2019, Litman2017, Pereira2017, Sheller2018}, and they can persist even in dockless systems. For example, in Shenzhen, China, both usage and perceived accessibility of dockless bicycles were higher in wealthier and more central neighborhoods~\cite{Zhou2024}. Reviews of cycling equity similarly report uneven distributions of infrastructure coverage, cycling activity, destination accessibility, and related health benefits, systematically favoring socially advantaged urban districts~\cite{Cunha2023}. Beyond physical access, informational barriers also matter: disadvantaged residents may be less aware of systems or targeted programs due to weaker user networks and limited access to smartphones or the internet~\cite{McNeil2018}. In Montreal, awareness of \textit{BIXI} increased mainly among higher-education groups and those living closer to stations, while lower-education groups remained less informed~\cite{Bernatchez2015}.

Encouragingly, evidence suggests that more equitable outcomes are achievable through inclusive design and policies. Interest in bike-sharing can be as high among disadvantaged groups as among advantaged ones~\cite{McNeil2018}, and per-capita usage can even be higher when barriers are reduced~\cite{Ogilvie2012}. In Washington, D.C., introducing \acp{e-b} alongside \acp{m-b} broadened the user base toward more racially diverse, lower-income, and more gender-balanced ridership~\cite{Bonilla-Alicea2020}. This implies that disparities are often driven less by preferences than by structural constraints, and that expanding coverage into underserved areas, maintaining affordability, and implementing targeted outreach can improve participation among marginalized communities~\cite{Goodman2014, McNeil2018, Ogilvie2012, Ursaki2015}.
